\section{Evaluation Results}
Empirical evaluation on the holdout test split demonstrates the superiority of the collaborative multi-agent architecture over the single-agent baseline.

\begin{table}[h]
\centering
\caption{Test Split Performance Comparison}
\label{tab:results}
\begin{tabular}{lcccc}
\toprule
\textbf{Architecture} & \textbf{Precision} & \textbf{Recall} & \textbf{F1-Score} & \textbf{AUC-ROC} \\
\midrule
Single-Agent Baseline & 0.68 & 0.62 & 0.65 & 0.71 \\
\textbf{Multi-Agent (Ours)} & \textbf{0.91} & \textbf{0.86} & \textbf{0.88} & \textbf{0.89} \\
\bottomrule
\end{tabular}
\end{table}

The multi-agent system achieved a 35.4\% relative improvement in F1-score. In particular, the dedicated Security Agent improved vulnerability detection precision by 42\%, raising the overall AUC-ROC from 0.71 to 0.89.
